En una experiència d'interferències sonores es fan servir dos altaveus, separats 4~m l'un de l'altre, que generen una ona acústica coherent i monocromàtica. Es localitzen diversos punts on la interferència és constructiva. Amb l'ajuda d'una cinta mètrica, es determina la distància de cadascun d'aquests punts als dos altaveus, concretament $R_1$ és la distància a un altaveu i $R_2$ és la distància a l'altre altaveu. Els resultats s'indiquen a la taula adjunta.

\figura[0.4]{fig1.png}

\begin{apartat}{1,25}
Ompliu la taula amb el càlcul del valor absolut de la diferència entre $R_1$ i $R_2$, i el número $n$ que identifica la línia d'interferència constructiva. Quin és el valor aproximat de la longitud d'ona i de la freqüència del so generat pels altaveus?
\end{apartat}

\begin{apartat}{1,25}
Quin és el nombre de punts nodals que apareixen al llarg de la línia que uneix els dos altaveus? Per què quan escolteu música amb un aparell estereofònic (dos altaveus) no es produeixen interferències?
\end{apartat}

\begin{center}
\begin{tabular}{|>{\raggedright}p{3cm}|c|c|>{\centering}p{2.2cm}|p{1.5cm}|}
\hline
Punt d'interferència constructiva & $R_1$ (m) & $R_2$ (m) & $|R_1-R_2|$ (m) & $n$ \\
\hline
1 & 2,54 & 2,51 & & \\ \hline
2 & 3,49 & 2,41 & & \\ \hline
3 & 1,82 & 2,89 & & \\ \hline
4 & 2,31 & 3,24 & & \\ \hline
5 & 4,58 & 2,63 & & \\ \hline
6 & 4,73 & 2,71 & & \\ \hline
7 & 5,27 & 2,36 & & \\ \hline
\end{tabular}
\end{center}

\dades{La velocitat del so a l'aire és de 343~m/s.}
